\section*{第 \ref{ch:b2:asexual-reproduction} 章——植物的无性生殖与克隆}

\begin{solution}{exo:b2:asexual-reproduction:1}
\omterm{def:b2:asexual-reproduction:clone}{克隆}：由一个亲本经
有丝分裂产生的遗传上相同的后代。\omterm{def:b2:asexual-reproduction:clone}{基株}：遗传上的个体，即由一个
受精卵长出的全部植株。\omterm{def:b2:asexual-reproduction:clone}{分株}：\omterm{def:b2:asexual-reproduction:clone}{基株}中一个生理上独立的单位。
器官：匍匐茎（草莓）、\omterm{prop:b2:asexual-reproduction:organs}{根状茎}（偃麦草）、\omterm{prop:b2:asexual-reproduction:organs}{块茎}
（马铃薯）、\omterm{prop:b2:asexual-reproduction:organs}{鳞茎}（洋葱）、根蘖（白杨）、叶上小植株
（\emph{Kalanchoe}）。
\end{solution}

\begin{solution}{exo:b2:asexual-reproduction:2}
接穗与\omterm{prop:b2:asexual-reproduction:horticulture}{砧木}的维管形成层必须对齐并
保持接触；愈伤组织跨越切口，并在其中分化出木质部和韧皮部。
\omterm{prop:b2:asexual-reproduction:horticulture}{砧木}提供根（水分、无机盐、活力、
耐土壤性、大小控制），接穗提供结果枝（即
品种）。两部分各自保留自己的基因组，细胞从不融合：
两种\omterm{def:b2:meiosis-heredity:vocabulary}{基因型}并列，构成嵌合体；杂种则会
在每个细胞中携带两种基因组的混合物。
\end{solution}

\begin{solution}{exo:b2:asexual-reproduction:3}
\omterm{def:b2:asexual-reproduction:clone}{营养增殖}：植物体的一部分生根，长成新植株。
\omterm{def:b2:asexual-reproduction:apomixis}{无融合生殖}：\omterm{def:b2:angiosperm-reproduction:seed}{种子}不经受精即形成，其胚是母体的\omterm{def:b2:asexual-reproduction:clone}{克隆}，
\omterm{def:b2:asexual-reproduction:apomixis}{孤雌生殖}：动物由未受精的卵发育而来。
只有\omterm{def:b2:asexual-reproduction:apomixis}{无融合生殖}保留了\omterm{def:b2:angiosperm-reproduction:seed}{种子}，包括\omterm{def:b2:angiosperm-reproduction:seed}{种皮}、\omterm{def:b2:angiosperm-reproduction:seed}{休眠}
与扩散。
\end{solution}

\begin{solution}{exo:b2:asexual-reproduction:4}
对外植体进行消毒，接种于含糖、无机盐、
维生素及适当平衡的生长素和细胞分裂素的培养基上（诱导愈伤组织）；
提高细胞分裂素与生长素的比例以诱导芽形成；
每隔几周继代培养芽；转入富含生长素的培养基生根；
小植株在高湿度下炼苗，然后移入土壤。
\end{solution}

\begin{solution}{exo:b2:asexual-reproduction:5}
每株植物被替换为 $1 + 24 = 25$：4 年后达到 $25^{4} =
\num{390625}$ 株，需要 \qty{39000}{m^2}，接近四公顷。
空间、光照和养分的限制，匍匐茎的延伸距离短，以及
拥挤下的植株死亡，使指数增长阶段在
两年左右停止；此后\omterm{def:b2:asexual-reproduction:clone}{克隆}以前沿的形式推进。
\end{solution}

\begin{solution}{exo:b2:asexual-reproduction:6}
$\num{14000}/130 \approx 110$ 代树干。圆盘面积为
\qty{4.3e5}{m^2} 时，其半径 $\sqrt{4.3\times 10^{5}/\pi} = \qty{370}{m}$；
每年推进 $370/\num{14000} = \qty{0.026}{m}$，
不足三厘米。
\end{solution}

\begin{solution}{exo:b2:asexual-reproduction:7}
$p_t = 2^{t}p_0/(1 + (2^{t} - 1)p_0)$：$t = 5$ 时为 $0.0032$；$t = 10$
时为 $0.1024/1.1023 = 0.093$；$t = 15$ 时为 $3.28/4.28 = 0.77$。若其适合度
低于有性谱系的一半，则永远不会扩散。相对适合度为 $w$ 时，
$q' = 2wq$，在 $w <
\tfrac12$ 时衰减。
\end{solution}

\begin{solution}{exo:b2:asexual-reproduction:8}
$10 : 1$：根；$1 : 1$：未分化的愈伤组织；$1 : 10$：芽。
插条已带有芽，需要生根，生长素可在切口基部
诱导根的形成。
\end{solution}

\begin{solution}{exo:b2:asexual-reproduction:9}
1、10、100、\num{1000}、\num{10000} 千克：第三季后仅有
\qty{1}{t}，第四季后达 \qty{10}{t}，共需四季。真实\omterm{def:b2:angiosperm-reproduction:seed}{种子}
每株可产数千粒\omterm{def:b2:angiosperm-reproduction:seed}{种子}，但马铃薯为\omterm{def:b2:meiosis-heredity:vocabulary}{杂合}四倍体，
因此实生苗会发生分离，没有一株能成为该品种。
\end{solution}

\begin{solution}{exo:b2:asexual-reproduction:10}
$e^{-5} = 0.0067$：大种群中约有 670 个无\omterm{def:b2:genome-diversification:mutation}{突变}个体，
小种群中约有 7 个。7 个个体可能在几代之内偶然地全部未能留下后代
（每代中它们都不繁殖的概率约为 $e^{-7}$），
此后最优类别永久消失，次优类别便成为最优：这是棘轮的一次咔嗒，
并会反复发生，
直至负荷不堪承受。
在 670 个个体的大种群中，损失几乎不可能发生——
除非瓶颈效应使 $N$ 缩小，或\omterm{def:b2:genome-diversification:mutation}{突变}率上升；
棘轮只会朝一个方向转动。
\end{solution}

\begin{solution}{exo:b2:asexual-reproduction:11}
在\omterm{def:b2:asexual-reproduction:clone}{克隆}中，每个宿主都易感，病原菌随孢子传播的速度
不受阻碍地扩散。在有性种群中，它所能克服的\omterm{def:b2:meiosis-heredity:vocabulary}{基因型}
只是 $2^{10} = 1024$ 种组合之一，
仅约占千分之一的植株，且散布在抗性邻株之间，
因此流行病无法蔓延。“Lumper”曾是在一个国家大部分地区
种植的单一\omterm{def:b2:asexual-reproduction:clone}{克隆}；晚疫病侵入时，每一株都毫无抗性，
而在引进其他品种之前，
没有可用于育种的变异。
\end{solution}

\begin{solution}{exo:b2:asexual-reproduction:12}
\omterm{prop:b2:asexual-reproduction:whysex}{穆勒棘轮}、有利\omterm{def:b2:genome-diversification:mutation}{突变}的组合以及红皇后效应，
各自都可使有性生殖在大型、受寄生、不断变化的世界中
获得超过两倍的益处。蛭形轮虫是个棘手的例子：
四千万年来没有雄性，也未出现棘轮式的崩溃。
它们似乎借助极端的抗干燥能力逃脱（这会杀死其寄生生物），
摄取外源 DNA，并
依靠异常高效的 DNA 修复——
这些例外恰好说明了通常规律所要求的条件。
\end{solution}

\begin{solution}{pb:b2:asexual-reproduction:1}
\textbf{1.} $4\times 10^{4}$ 株。
\textbf{2.} $1 + 6\times 2 = 13$。
\textbf{3.} \num{1300}；\num{16900}；\num{219700}。
\textbf{4.} $13^{t} = 400$：$t = \ln 400/\ln 13 = 5.99/2.56 = 2.3$
年——第三年内即被占满。
\textbf{5.} 一公顷圆盘的半径 $\sqrt{10^{4}/\pi} = \qty{56}{m}$；
以每年 \qty{0.5}{m} 的速度，需 \num{113} 年。
\textbf{6.} 一百个分散的奠基者以指数方式从一百道前沿
在两年内铺满田地；若只有一个奠基者，它
在两年内即可占满其邻近区域，此后再缓慢蔓延一个世纪。
\textbf{7.} $5^{n} \ge 10^{6}$：$n \ge 8.6$，需九次继代
（$1.95\times 10^{6}$ 个芽），36~周。
\textbf{8.} $0.8\times 1.95\times 10^{6} = 1.6\times 10^{6}$ 株。
\textbf{9.} $0.98$；$0.98^{10} = 0.82$。
\textbf{10.} $10\times 0.98 = 9.8$ 个无病毒品系（预期值）；
由十个插条中得到一个。
\textbf{11.} 病毒经胞间连丝和韧皮部运输，
而茎尖圆顶没有维管连接，且分裂速度快于
病毒的扩散速度，因此最外层十分之一毫米
通常未被感染。
\textbf{12.} 实验室方法：可立即获得整齐、无病毒的植株，但
成本高昂，且培养物中的任何\omterm{def:b2:genome-diversification:mutation}{突变}或污染都会
被放大一百万倍。匍匐茎：免费且就地可得，但
会损失三个季节，且母株中的每种病毒都会被一并带走。
\textbf{13.} \omterm{def:b2:asexual-reproduction:apomixis}{无融合生殖}个体数为 $pN$，每个产 $F$ 粒\omterm{def:b2:angiosperm-reproduction:seed}{种子}；有性个体
$(1 - p)N$，每个产 $F/2$ 粒\omterm{def:b2:angiosperm-reproduction:seed}{种子}（其一半投入是\omterm{def:b2:plant-life-cycles:heterospory}{花粉}，\omterm{def:b2:plant-life-cycles:heterospory}{花粉}
不产生\omterm{def:b2:angiosperm-reproduction:seed}{种子}）；\omterm{def:b2:asexual-reproduction:apomixis}{无融合生殖}个体在\omterm{def:b2:angiosperm-reproduction:seed}{种子}中的占比为 $Fp/(Fp + F(1 - p)/2)
= 2p/(1 + p)$。
\textbf{14.} $p_5 = 0.032/1.031 = 0.031$；$p_{10} = 1.024/2.023 =
0.51$：第 10 代时超过一半。
\textbf{15.} $p' = 2p(1 - cp)/\bigl(2p(1 - cp) + (1 - p)\bigr)$。
\textbf{16.} $p' = p$ 在 $2(1 - cp^*) = 1$ 时成立：$p^* = 1/2c = 0.625$。
在 $p^*$ 以下时 $2(1 - cp) > 1$，$p$ 上升；以上则下降：该平衡是稳定的。
\textbf{17.} $p^* \ge 1$ 在 $c \le \tfrac12$ 时成立：\omterm{def:b2:asexual-reproduction:apomixis}{无融合生殖}个体取得优势，
除非寄生生物在其普遍存在时使其损失超过产出的一半。
\textbf{18.} $c = 0.6$：$p^* = 0.83$。要使有性生殖保留在种群中，
寄生生物在\omterm{def:b2:asexual-reproduction:apomixis}{无融合生殖}个体频率较高时，必须使\omterm{def:b2:asexual-reproduction:clone}{克隆}的
适合度降幅超过两倍的优势——红皇后必须
快速奔跑。
\textbf{19.} 三套染色体在\omterm{def:b2:meiosis-heredity:meiosis}{减数分裂}中无法配对，因此其
配子不平衡；它被排除在重组之外，同时
\omterm{prop:b2:asexual-reproduction:whysex}{穆勒棘轮}与寄生生物对其均无阻碍地起作用——这正是
为什么\omterm{def:b2:asexual-reproduction:apomixis}{无融合生殖}的蒲公英谱系年轻，并且不断由
有性祖先重新产生。
\textbf{20.} $10^{-9}\times 4.8\times 10^{8} = 0.48$，每次分裂。
\textbf{21.} $2\times \num{14000}\times 20 = 5.6\times 10^{5}$
次分裂；$0.48\times 5.6\times 10^{5} = 2.7\times 10^{5}$
个\omterm{def:b2:genome-diversification:mutation}{突变}。
\textbf{22.} $2.7\times 10^{5}/4.8\times 10^{8} = 5.6\times 10^{-4}$，即
占基因组的比例，约 \num{5400} 个编码序列差异。
\textbf{23.} 分生组织中细胞谱系间的竞争（分裂较慢的\omterm{def:b2:genome-diversification:mutation}{突变}系被取代），
以及\omterm{def:b2:asexual-reproduction:clone}{分株}间的竞争（携带有害\omterm{def:b2:genome-diversification:mutation}{突变}的树干
产生的根蘖较少）。这种选择较弱，因为大多数体细胞\omterm{def:b2:genome-diversification:mutation}{突变}是\omterm{def:b2:meiosis-heredity:vocabulary}{杂合}且被掩盖的，
且没有任何机制使其暴露：
没有\omterm{def:b2:meiosis-heredity:meiosis}{减数分裂}使其\omterm{def:b2:meiosis-heredity:vocabulary}{纯合}，也没有单细胞受精卵阶段
可在一个细胞中检验整个基因组。
\textbf{24.} 并不一致：其树干在数十万
个位点上存在差异。但这些差异大多是中性且被掩盖的，因此
该\omterm{def:b2:asexual-reproduction:clone}{克隆}实际上是一个\omterm{def:b2:meiosis-heredity:vocabulary}{基因型}，
即同一\omterm{def:b2:asexual-reproduction:clone}{基株}体细胞变异的镶嵌体。
\textbf{25.} 2.3 年后田地被占满；一个分生组织
在九个月内可产生 $1.6\times 10^{6}$ 株植物；当 $c = 0.8$ 时，\omterm{def:b2:asexual-reproduction:apomixis}{无融合生殖}个体
稳定于 $p^* = 0.625$，有性个体保留 \qty{37.5}{\%}。
\end{solution}
